\begin{recipe}
[% 
    preparationtime = {\unit[30]{min}},
    portion = {\portion{4}},
    %source = {},
    calory = {695kcal | 18g Protein | 103g KH | 22g Fett},
    price = {2,10€},
    created = {17.05.2026},
    rating = {3},
]
{Pasta nach Art Puttanesca}

\graph{% pictures
    big=pic/hauptgericht/16_pasta-puttanesca
}

\introduction{%
Eine schnelle, intensive Tomatenpasta mit Oliven, Kapern, Knoblauch und Chili. Salzig, fruchtig, leicht scharf und genau die Art Gericht, die nach mehr Aufwand schmeckt als sie eigentlich ist.
}

\ingredients{
    500 g & Penne\index{Getreide!Pasta}\\
    60 g & Zwiebeln\index{Gemüse!Zwiebel}\\
    2 Zehen & Knoblauch\index{Gemüse!Knoblauch}\\
    30 g & Kapern\index{Gemüse!Kapern}\\
    100 g & Grüne Oliven, entsteint\index{Obst!Olive}\\
    100 g & Schwarze Oliven, entsteint\index{Obst!Olive}\\
    1 & Frische Chilischote\\
    3 EL & Tomatenmark\\
    2 EL & Getrocknete Tomaten\index{Gemüse!Tomate}\\
    50 ml & Olivenöl\\
    500 g & Geschälte Tomaten / Passata\index{Gemüse!Tomate}\\
    1 Spitze & Meersalz\\
    & Fruchtiges Olivenöl zum Servieren
}

\preparation{%
    \step%
    Einen großen Topf mit Wasser aufsetzen und kräftig salzen. Die Pasta darin sehr al dente kochen, da sie später noch kurz in der Sauce fertig gart.
    
    \step%
    Zwiebeln und Knoblauch schälen und fein würfeln. Kapern grob hacken. Oliven ebenfalls grob schneiden. Chilischote entkernen und fein schneiden.
    
    \step%
    Olivenöl in einem großen Topf erhitzen. Zwiebeln darin glasig anschwitzen.
    
    \step%
    Knoblauch, Chili, Kapern, Oliven und getrocknete Tomaten zugeben und etwa 10 Minuten bei mittlerer Hitze farblos dünsten.
    
    \step%
    Tomatenmark einrühren und weitere 10 Minuten langsam mitrösten.
    
    \step%
    Geschälte Tomaten oder Passata zugeben und die Sauce mit einem Stabmixer kurz anpürieren. Sie soll nicht komplett glatt werden, sondern noch etwas Struktur behalten.
    
    \step%
    Sauce etwa 10 Minuten leicht köcheln lassen.
    
    \step%
    Pasta abgießen und dabei etwas Kochwasser auffangen. Nudeln in die Sauce geben und mit etwas Kochwasser etwa 1 Minute fertig garen.
    
    \step%
    Petersilie nach Geschmack zugeben, mit Zitronenabrieb, Pfeffer und etwas Olivenöl abschmecken.
}

\hint{
    Die Sauce ist durch Oliven und Kapern schon ziemlich salzig. Also beim finalen Abschmecken lieber vorsichtig sein und nicht direkt Salz reinballern.
}

\end{recipe}